\section{Introduction}
\label{sec:intro}

Dense retrieval underpins many modern systems, including Retrieval-Augmented
Generation (RAG)~\citep{lewis2020retrieval}, question answering, and
semantic search~\citep{karpukhin2020dense, reimers2019sbert}. A bi-encoder
maps queries and passages to vectors, and relevance is scored by vector
similarity. Long documents are split into passages, so the indexed unit is a
passage that is encoded once and stored. These corpus vectors must reside in an
index in Random Access Memory (RAM) for low-latency
search~\citep{johnson2021billion, malkov2018efficient}. This memory is the dominant cost at scale. A single
BGE-base embedding is a 768-dimensional float32 vector of 3{,}072 bytes, so a
corpus of one billion passages needs roughly 2.9~TB of RAM for the embeddings
alone. Memory-optimized cloud instances are billed accordingly, so reducing
the index size translates directly into lower cost~\citep{aws2024pricing}.

Embedding compression addresses this cost, and two broad families exist.
\emph{Post-hoc} methods compress a fixed embedding after training: scalar
quantization to INT8, binary quantization with Hamming
search~\citep{shakir2024embedding}, Product Quantization
(PQ)~\citep{jegou2010product}, and recent oblivious quantizers such as
TurboQuant~\citep{zandieh2025turboquant}. \emph{Training-time} methods bake
the compression into the objective: binarization-aware training through a
Straight-Through Estimator (STE)~\citep{bengio2013estimating, courbariaux2016bnn}
or an annealed tanh surrogate~\citep{yamada2021efficient}, and Matryoshka
Representation Learning (MRL)~\citep{kusupati2022matryoshka}, which front-loads
information so that truncated prefixes remain usable.

\paragraph{The gap.}
Despite a rich set of methods, a practitioner cannot easily decide which to
use. The reason is methodological: post-hoc and training-time methods are
studied on different encoders, with different training data and benchmarks, and
they rarely meet on common ground. The broadest post-hoc
comparison~\citep{correct2025} does vary the backbone, but omits training-time
methods, while the training-time studies~\citep{qama2025, yamada2021efficient}
fix a single backbone or a narrow pair. So it is unknown whether the best method
depends on the encoder's retrieval prior, and whether training-time compression
is worth its added cost relative to a simple post-hoc transform.

\paragraph{This work.}
We run a controlled comparison. Three encoders with deliberately different
retrieval priors, BGE~\citep{xiao2023c}, RoBERTa~\citep{liu2019roberta}, and
MPNet~\citep{song2020mpnet}, are fine-tuned under identical conditions on the
same data. Each is then compressed by every method above and by stacked
combinations, and all are evaluated on NanoBEIR~\citep{nanobeir2024} with
paired significance testing. Because the only variable that changes is the
compression method (and the backbone), differences are attributable to the
method rather than to confounds of data or benchmark.

\paragraph{Findings.}
The headline result is that no single method dominates, and the winner depends
on the encoder. For BGE, which is pre-trained at scale with a contrastive
retrieval objective, post-hoc PQ leads the Pareto frontier at every tested
compression ratio. For RoBERTa and MPNet, which lack retrieval pre-training,
training-time binarization leads instead. Most surprisingly, binarization-aware
training \emph{exceeds the full-precision baseline} on RoBERTa, raising Mean Reciprocal Rank at 10 (MRR@10)
from 0.530 to 0.590, even though binarization removes information.

We explain this with the geometry of the embedding space. Encoders without
retrieval pre-training produce anisotropic embeddings that occupy a narrow
cone~\citep{ethayarajh2019contextual}, which saturates similarities and harms
fine-grained ranking. We quantify this with the \emph{effective dimension}, or
participation ratio~\citep{litwinkumar2017optimal, gao2017theory}.
Binarization-aware training spreads variance across more axes, raising the
participation ratio from 11\% to 30\% of 768 dimensions for RoBERTa. The gain
shows up in ranking, not coverage: MRR@10 improves while Recall@10 remains
unchanged.
The effect is present from the first checkpoint and does not close with more
training, so it is structural rather than a symptom of under-training. Because
ranking under this metric depends on angular separation rather than per-axis
precision, one bit per dimension is sufficient to realize the benefit.

\paragraph{Contributions.}
\begin{itemize}
  \item A controlled, cross-encoder comparison of post-hoc, training-time, and
        stacked embedding compression for dense retrieval, reported as Pareto
        frontiers with paired significance tests
        (Sections~\ref{sec:method} and~\ref{sec:results}).
  \item The finding that binarization-aware training can exceed full precision
        on encoders without retrieval pre-training, with a mechanism:
        binarization raises the effective dimension of the embedding space and
        improves ranking without changing recall
        (Section~\ref{sec:results}, Section~\ref{sec:discussion}).
  \item Evidence that the best gradient surrogate and the value of stacking are
        backbone-dependent, including that annealed tanh combines with MRL more
        cleanly than STE, and that MRL$+$PQ extends the frontier to high
        compression (Section~\ref{sec:results}).
  \item A practical method-selection guide indexed by encoder type and memory
        budget (Section~\ref{sec:discussion}), released together with all code
        and an interactive evaluation dashboard for browsing the full results.

\end{itemize}
